\section{LLM Usage \& Agent Improvement Process}

\subsection{LLM Models Used}
\begin{itemize}
    \item \textbf{Big Pickle:} Evaluated via the OpenCode Zen platform, this
    model is specifically optimized for automated coding agents. Featuring an
    extensive context window of approximately 200k tokens, it serves as an
    ideal tool for rapid initial prototyping. However, its primary limitation
    lies in the absence of advanced reasoning features.

    \item \textbf{DeepSeek V4 Flash Free:} Accessed through the OpenCode
    platform, this model was deployed to facilitate technical drafting,
    content refinement, and structural editing, such as database design
    documentation. It demonstrates high efficiency in generating structured
    explanations and concise academic prose, though it necessitates human
    oversight to ensure rigorous factual accuracy and alignment with project
    specifications.
\end{itemize}

\subsection{Evaluation and Improvement Cycle}
Explain how agent output was evaluated after each Phase 2 task. Summarize the
problems found, the evidence used to identify them, and the corresponding
changes to prompts, \texttt{AGENTS.md}, skills, validation steps, or review
checklists.

\begin{table}[htbp]
    \centering
    \begin{tabularx}{\textwidth}{p{0.16\textwidth}X X X}
        \toprule
        \textbf{Task} & \textbf{Evaluation evidence} & \textbf{Issue found} & \textbf{Improvement made} \\
        \midrule
        % TODO: Add rows for deliverables 08--16.
        08--16 & To be completed & To be completed & To be completed \\
        \bottomrule
    \end{tabularx}
    \caption{Phase 2 agent evaluation and improvement log}
    \label{tab:phase2-agent-improvement}
\end{table}